\section[谱理论与层理论]{谱理论与层理论（1 月 26 日）}
\srcpage{8}

今天我们开始讲得更慢、更深；接下来几讲，我们会聚焦于棱镜这幅图景：自守形式从左边进入，被谱分解成 Galois 表示。不过，在进入谱分解之前，先问一句：什么是谱？

我们从几何谈起。几何对不同的人可能确切地意味着不同的东西，但几何中有某种「空间」的概念；给定一个空间，就有一个函数构成的交换环，乘法逐点取。这是代数几何的基本出发点：要研究空间，就研究其函数代数。事实上，我们可以把取函数环看成从空间到交换环的反变函子 $\mathcal{O}$：给定空间之间的映射 $X\to Y$，我们想要函数的拉回映射，它保持逐点加法与乘法。

谱的基本想法是反其道而行：从交换代数出发，用它造出一个空间。范畴论精确地澄清了我们要做什么：应当有一个从交换代数到空间的谱函子 $\Spec$，它满足一个万有性质，而这个性质可简洁地表述为：它是 $\mathcal{O}$ 的右伴随。也许为了完全精确，我们得说清关心哪类空间与函数；这里有不同选项，但在这些不同情形中，你都有这个基本问题。

明确地说，$\Spec$ 是 $\mathcal{O}$ 的右伴随，意思是：对每个交换代数 $R$，空间映射 $X\to\Spec R$ 的集合自然同构于 $\Hom_{\mathrm{Ring}}(R,\mathcal{O}(X))$。你可以在函子范畴 $\mathrm{Fun}(\mathrm{Spaces},\mathrm{Set})$ 中造出一个「弱解」：$\Spec R$ 把 $X$ 送到 $\Hom_{\mathrm{Ring}}(R,\mathcal{O}(X))$；然后可以问它是不是真解，即是否被某个真正的空间表示。我们还想要 $\mathcal{O}(\Spec R)=R$，这里的「$=$」指「自然同构」。

有几种不同的设定，在其中这套做法相当奏效。

\srcpage{9}

\begin{example}
设「空间」指「有限集」，「函数」指取值于 $k$ 的函数，$k$ 是你最喜欢的域。（我最喜欢的域是 $\mathbb{C}$。你呢？）此时 $\mathcal{O}(X)$ 是函数环 $\{X\to k\}$，即 $\prod_{x\in X}k$，也就是对角 $|X|\times|X|$ 矩阵构成的代数。所以在这样的设定里，有许多代数并不是空间上的 $k$-函数代数，例如 $k[x]/(x^2)$。
\end{example}

\begin{example}[Gelfand]
设「空间」指局部紧 Hausdorff 空间，「函数」指取值于 $\mathbb{C}$、在无穷远处消失的连续函数：$\mathcal{O}(X):=C_0(X)$。这具有交换 $C^*$-代数的结构，$*$-运算由复共轭给出。\footnote{关于这个对偶的一些讨论，参看 [aHRW09]。}
\end{example}

在这样的设定中，有一个很好的 $\Spec$ 函子，从交换 $C^*$-代数到局部紧 Hausdorff 空间，即 Gelfand 谱。给定这样的代数 $A$，$\operatorname{mSpec}A$ 定义为 $A$ 的极大理想构成的空间。这些极大理想与集合 $\Hom_{C^*}(A,\mathbb{C})$ 等同，而后者可以很有益地想成 $A$ 的「点」：$\mathbb{C}$ 是一个点上的函数，所以从启发式上说，这就是映射 $\mathrm{pt}\to\operatorname{mSpec}A$。换个说法，这些是 $A$ 的一维酉表示。

\begin{theorem}[Gelfand–Naimark]
$(\mathcal{O},\operatorname{mSpec})$ 是一对反变范畴等价，从局部紧 Hausdorff 空间到交换 $C^*$-代数。
\end{theorem}

于是连续拓扑的世界完全可以由代数知晓。

在这个等价下，紧 Hausdorff 空间（即好的那些）与含幺交换 $C^*$-代数（即好的那些）互换。这是因为取值为 $1$ 的常函数在 $C^*$ 范数下有界，当且仅当底空间紧。

\begin{example}
再粗一点儿的例子：设「空间」指测度空间，「函数」指 $L^\infty(X)$，这就落入了交换 von Neumann 代数的世界。要把一切说精确，这里还该多讲几句。同样有反变范畴等价；而这一次对象的同构类并不多：点的有限并、点的可数并、区间，以及区间与一些点的并。
\end{example}

\begin{example}
代数几何是研究得最充分的例子，但它不像其他一些例子那样奏效得那么好。此时空间指局部环化空间，即拓扑空间 $X$ 连同环层 $\mathcal{O}_X$，后者带有我们此处不打算深究的性质；环则指交换环。取 $\mathcal{O}_X$ 的整体截面，定义了到交换环的反变函子，而它有右伴随 $\Spec$，即环的谱。它并非本质满：像中的东西称为仿射概形；在仿射概形构成的全子范畴上，$(\mathcal{O},\Spec)$ 是一对反变范畴等价。
\end{example}

不幸的是，这没有涵盖许多重要的例子：局部看起来像仿射概形的局部环化空间。代数几何中许多重要对象，例如射影直线，都是以这种方式由仿射概形搭建出来的，但它们本身不是仿射的。于是你把几何的概念稍微扩大一点。但在这个视角下，确实有一些有用的东西，它们是函子到集合之表示问题的弱解而非强解。既然如此，为什么不干脆就用这些弱解来做几何呢？

\begin{example}[Quillen–Sullivan 有理同伦论]
在同伦论中，你可能想研究同伦范畴，即把（局部紧弱 Hausdorff）拓扑空间中的同伦等价倒过来得到的范畴。\footnote{做这件事要小心，以避免集合论问题，但这是可以做到的。}这里「环」指分次交换环，函子是 $H^*(-;\mathbb{Z})$。上同调很好，但它不够灵活，无法建立好的谱理论。
\end{example}

若记住整个（局部紧弱 Hausdorff）拓扑空间范畴，并让函子取有理上链，则可以得到更好的对应；有理上链取值于交换微分分次 $\mathbb{Q}$-代数（CDGA）范畴。Quillen–Sullivan 证明了这定义了一套很好的谱理论：限制到单连通空间，它给出从「模去有理同伦等价的空间」到「单连通 $\mathbb{Q}$-CDGA」的范畴等价。\footnote{可以稍作推广到幂零空间，但必须对 $\pi_1$ 加某种条件。这有点像代数几何中的非仿射性，不过实践中它的行为略有不同。}

Mandell 有 $p$ 进版本的类似陈述；Allen Yuan [Yua19] 最近有整版本的陈述，只不过用的是带有稍多结构的上链。

\begin{remark}
数学中「谱」的大多数用法——代数几何的、算子理论的、甚至数学物理味的——都彼此相关。唯一的例外是同伦论学家所说的谱，那意味着不同的东西。要提防这个常见的混淆来源。本课程中，「谱」大多按整的意义（而非有理的意义）来用。
\end{remark}

\srcpage{10}

回到谱分解。设交换环 $R$（常常是 $k$-代数；对我们来说总是 $k$-代数，其中 $k=\mathbb{C}$）作用在一个模（对我们来说是复向量空间）$V$ 上。这是代数同态 $R\to\End(V)$ 的数据。我们对谱分解的看法是：想把这个模层化、局部化、铺展开，或者在 $\Spec R$ 上把它谱分解。为此我们要用到 $\mathrm{Mod}_R$ 是对称幺半的：有张量积 $\otimes_R:\mathrm{Mod}_R\times\mathrm{Mod}_R\to\mathrm{Mod}_R$。于是可以定义与模 $V$ 相伴的层 $\ul{V}$ 为
\[
\ul{V}(U) := V\otimes_R\mathcal{O}(U).
\]
一个直接的推论是：你可以谈元素 $v\in V$ 的支撑，即子集 $\operatorname{supp}(v)\subset X$。

例如，若 $X$ 是有限集，于是 $R=\prod_{x\in X}\mathbb{C}$，则
\[
V=\bigoplus_{x\in X}V_x.
\]
由这个构造产生的层，值得有个名字。

\begin{definition}
$\Spec R$ 上的拟凝聚层（quasicoherent sheaf），就是以这种方式由 $R$-模得到的层。若 $R$-模是有限生成的，该层称为凝聚层（coherent sheaf）。
\end{definition}

并非所有概形都是仿射的，所以我们说（拟）凝聚层是那些局部为此形式的层。

（TODO：我漏掉了一个例子，其中谱分解就是向量空间通常按特征空间的分解。这是指 $\mathbb{C}[x]$ 中 $x$ 按所论矩阵作用吗？还是更像有限集那个例子？这个例子有点太基本：点是既开又闭的，所以特征空间既是 hom 又是张量。一般情形并非如此。）

那么，谱定理在这个语境中如何出现？设 $V$ 是向量空间（不必有限维，虽然一般需要某种好的拓扑），$M\in\End(V)$。我们把这个矩阵看成代数映射：$\Hom_{\mathrm{Set}}(\mathrm{pt},\End(V))$。现在 $\End(V)$ 有许多额外结构——它是结合 $\mathbb{C}$-代数，虽然不交换。所以应当有一个伴随
\[
\Hom_{\mathrm{Set}}(\mathrm{pt},\End(V))=\Hom_{\mathrm{Alg}_{\mathbb{C}}}(\mathrm{Free},\End(V)),
\]
其中 $\mathrm{Free}$ 是一个生成元上的自由结合代数，即 $\mathbb{C}[x]$。所以与其想矩阵 $M$，我们将想「$V$ 是 $\mathbb{C}[x]$-模」这个数据。在某些语境中，这称为「函数演算」（functional calculus）——一旦有了矩阵，你就可以用这个矩阵的多项式来作用。我们这里的做法是把所有这些放在一起想，而不只是想 $M$。事实上，若拓扑足够好，你可以完备化，从而理解 $\mathbb{R}$ 上的函数之类的东西，而不只是多项式。

无论如何，想法是：$M\in\End(V)$ 等价于 $V$ 是仿射直线 $\mathbb{A}^1:=\Spec\mathbb{C}[x]$ 上某个拟凝聚层的整体截面。这将是我们谱分解的基本情形：一族交换算子中最简单的情形，就是单个算子。

我们想研究 $V$ 如何在 $\mathbb{A}^1$ 上铺展开。利用主理想整环上模的结构理论，有短正合序列
\[
0\longrightarrow V_{\mathrm{tors}}\longrightarrow V\longrightarrow V_{\mathrm{free}}\longrightarrow 0,
\]
事实上它分裂。所以 $V_{\mathrm{free}}\cong\mathbb{C}[x]^{\oplus r}$，而挠部分直和分解为
\[
V_{\mathrm{tors}}\cong\bigoplus_{\lambda\in\Spec(V)}\widehat{V}_\lambda.
\]
自由部分我们称为连续谱，挠部分称为离散谱。具体地说，若 $V$ 支撑在 $\lambda\in\mathbb{A}^1$ 处，这就是说 $\lambda$ 是一个广义特征值，而 $\widehat{V}_\lambda$ 是（等价于）$M$ 的 Jordan 块。

特征向量 $Mv=xv=\lambda v$ 等价于 $\Hom_{\mathbb{C}[x]}(\mathbb{C}_\lambda,V)$ 的元素，这涉及的是商而不是子（TODO：这里漏掉了点什么）。而若你看连续谱，那里没有特征向量；也就是说，在 $\mathbb{A}^1$ 上，自由部分看起来像平凡丛的截面，而离散谱是各点处的摩天楼层。两种情形下我们都可以取纤维（商）。（TODO：我该画一张这样的图。）

\srcpage{11}

\begin{remark}
谱字典有许多版本。若谈 von Neumann 代数与可测空间，对应的谱定理是 von Neumann 谱定理：$M=A$ 是 Hilbert 空间 $V$ 上的自伴算子。把这个定理层化地重述，就是说有一个「层」，即 $\mathbb{R}$ 上的投影值测度（$\mathbb{R}$ 是谱），而算子 $A$ 可以重构为直积分
\[
A=\int_{\mathbb{R}}x\,\mathrm{d}\pi_A,
\]
这是关于投影值测度 $\pi_A$ 的积分。
\end{remark}

为了理解这句话，我们来解释什么是投影值测度。它是从每个可测子集 $U\subset\mathbb{R}$ 到 $V$ 上投影算子 $\pi_A(U)$ 的指派。不同的投影应当交换。所以这是某种 Hilbert 空间的层，只是意义特别弱。你可以把 $U$ 上的截面空间想成 $\operatorname{Im}(\pi_A(U))$，并且有可数可加性：
\[
U\longmapsto\langle w,\pi_A(U)(v)\rangle
\]
必须是 $\mathbb{R}$ 上的 $\mathbb{C}$-值测度。$A$ 的谱是 $\pi_A$ 的支撑。若 $V$ 有限维，直积分就变成直和，给出 $V$ 按 $A$-特征空间的分解。一般地，直积分会把连续谱也考虑进去。

在我们讨论过的其他谱设定中，也有类似图景的版本，例如 $C^*$-代数的 Hilbert $C^*$-模，等等。甚至还有同伦版本：若 $R=C^*(X)$ 是空间 $X$ 的，则 $\mathrm{Mod}_R$ 嵌入 $\operatorname{Loc}(X)$，即 $X$ 上局部常值（复形）层的范畴；而适当意义下有限生成的 $R$-模，对应 $X$ 上的一个（复形）层，并且以同样的方式有限生成。这与下述事实有关：$C^*(X)=\End(k_X)$，其中 $k_X$ 是常局部系统，且 $\ul{M}(U)=M\otimes_{C^*(X)}C^*(X)$。这里还有一些同伦细节要补，但一切都可以说精确；要点是这里也有故事的类比物：你有与空间 $\Spec R$ 相伴的代数 $R$，而 $R$-模铺展在 $\Spec R$ 上。

\begin{center}
$\ast\quad\ast\quad\ast$
\end{center}

谱分解在代数、空间与模的例子很多的地方最出彩。为此，我们转向 Fourier 理论。

在这个设定中，我们希望分布是 $V$ 的线性对偶 $\Hom(V,\mathbb{C})$。这里有许多弱特征向量，例如 $\delta$-函数，它们可能并不是真正的特征向量（例如前面看到的连续谱问题）。举例来说，若 $V=L^2(\mathbb{R})$，算子 $M=x$，你在 $L^2$ 世界里没法理解「支撑在 $x$ 处的函数」。但对分布可以这样做。对偶地，若 $M^\vee=\frac{\mathrm{d}}{\mathrm{d}x}$，我们有关于 $\lambda$ 的自然特征向量 $e^{i\lambda x}$，但它们不是 $L^2$ 的：它们生活在更大的空间里，由另一种范数控制。连续谱常常就是这副样子：你有一个直积分，它与直和不同。你所积分的东西其实并不是子集。例如，$\mathbb{A}^1$ 上的函数是一个点上的函数直接拼起来的，但 $\mathbb{A}^1$ 中一个点处的函数并不是 $\mathbb{A}^1$ 上函数的子集，而是商。没有挠的时候情况会好一些，但那会掩盖一般的故事。而这个一般的故事并不是奇怪的分析事实，因为它在代数中对于多项式也会出现。

下次我们会从这个视角谈 Fourier 理论，或者说阿贝尔对偶：把它当作交换算子的一个自然来源。阿贝尔群 $G$ 作用在向量空间 $V$ 上，是寻找大交换算子代数的好地方。我们会用这些算子把 $V$ 谱分解。本课程的目标——某种意义上也是 Langlands 纲领更广阔的目标——是这件事的非阿贝尔推广；而我们会花几周时间在阿贝尔的故事上。
